\documentclass[aps,pre,preprint,notitlepage,nofootinbib,superscriptaddress]{revtex4-2}

\usepackage[T1]{fontenc}
\usepackage{lmodern}
\usepackage{microtype}
\usepackage{amsmath,amssymb,mathtools,bm,amsthm}
\usepackage{mathrsfs}
\usepackage{booktabs,array}
\usepackage{xcolor}
\usepackage{hyperref}
\usepackage[nameinlink,noabbrev]{cleveref}
\usepackage{enumitem}

\hypersetup{
  colorlinks=true,
  linkcolor=blue!55!black,
  citecolor=blue!55!black,
  urlcolor=blue!55!black,
  pdftitle={Energy-Localization Closures for a Vortex-Filament Model of the Electron Mass},
  pdfauthor={Omar Iskandarani}
}

\newcommand{\dd}{\mathrm d}
\newcommand{\lC}{\bar\lambda_C}
\newcommand{\re}{r_e}
\newcommand{\Rop}{\operatorname{Rop}}
\newcommand{\calL}{\mathcal L}
\newcommand{\rhohorn}{\rho_{\mathrm{horn}}^{\mathrm{eff}}}
\newcommand{\rhofil}{\rho_{\mathrm{fil}}^{\mathrm{eff}}}
\newcommand{\rhoCeff}{\rho_{\mathrm C}^{\mathrm{eff}}}
\newcommand{\Gammaeff}{\Gamma_{\mathrm{eff}}}
\newcommand{\GammaC}{\Gamma_{\mathrm C}}

\newtheorem{proposition}{Proposition}
\newtheorem{remark}{Remark}
\newtheorem{definition}{Definition}


\begin{document}

\title{\texorpdfstring{Conditional Energy-Normalization Closures for a Vortex-Filament Electron Model\\\large Exact algebraic ratios, \(\alpha\)-power counting, and open constitutive gates}{Conditional Energy-Normalization Closures for a Vortex-Filament Electron Model}}

\author{Omar Iskandarani}
\email{info@omariskandarani.com}
\thanks{ORCID: 0009-0006-1686-3961}
\affiliation{Independent Researcher, Groningen, The Netherlands}
\date{2 August 2026; manuscript version 0.2.0}

\begin{abstract}
We audit three conditional energy-normalization closures used in a vortex-filament bookkeeping of the electron rest energy.  The first is the calibrated horn-envelope normalization
\(m_ec^2=2\pi\rhohorn v^2r^3\); the second is a declared leading line-energy closure
\(m_ec^2=\rhofil\Gamma^2\ell/3\); the third replaces the effective circulation by the electron-anchored Compton circulation scale \(\GammaC=h/m_e\).
For the fixed coefficients \(2\pi\) and \(1/3\), and with \(\Gamma=2\pi rv\), the first two closures obey the exact conditional identity
\[
\frac{\rhohorn}{\rhofil}
=\frac{2\pi}{3}\frac{\ell}{r}.
\]
On the calibrated chain \(r=(\alpha/2)\bar\lambda_C\) and on the candidate length
\(\ell=(3/8\pi)\bar\lambda_C\), this becomes \(1/(2\alpha)=68.518\).
This factor is a ratio of normalization supports defined by the two closures; it is not a theorem about the geometric volume of a sphere and a circular tube, nor does it identify either coefficient with a material, background, or resolved-core density.

The geometric candidate
\(\alpha_{\rm geom}^{-1}=(8\pi/3)\mathcal L_D^{(3_1)}\), the associated comparison diameter \(D_{\rm geom}\), and the required radius-normalized aspect ratio encode one and the same residual, \(-855.13\) ppm in \(\alpha\).
Replacing \(\Gammaeff\) by \(\GammaC=h/m_e\) produces the \(\alpha\)-free but electron-anchored density scale
\(\rhoCeff=(2/\pi)m_e/\bar\lambda_C^3\).
Neither \(\GammaC\) nor \(\rhoCeff\) is derived as a property of the substrate.
The ratio \(\Gammaeff/\GammaC=\alpha^2/4\) is therefore a calibrated comparison, not yet a dynamical circulation deficit.
A formal line-vortex interpretation of \(\GammaC\) at sub-Compton radii gives flow speeds exceeding \(c\); this conflicts with a universal-\(c\) relativistic interpretation of that flow, but not with the mathematics of a Galilean incompressible model by itself.

The analysis isolates five independent gates: dynamical selection of the ideal ropelength, computation of the line-energy coefficient, constitutive identification of the relevant density, derivation of the circulation sector, and Maxwell/source matching including the electromagnetic self-energy shape factor.
The model derives neither \(\alpha\), \(m_e\), nor a medium density.
Its present content is a transparent algebraic scaffold and a preregistered falsification programme.
\end{abstract}

\maketitle

\section{Introduction}
\label{sec:intro}

The proposal that particles may be represented by knotted or linked vortical structures has a long history \cite{Thomson1867}.  Quantized vortices are established excitations of superfluids \cite{Onsager1949,Feynman1955,Vinen1961,Donnelly1991}; knotted vortices have been generated and tracked in ordinary fluids \cite{KlecknerIrvine2013}; and tight knots have been used as organizing structures in phenomenological flux-tube spectra \cite{BuniyKephart2005,BuniyEtAl2014}.  Magnetic and hydrodynamic tube models likewise motivate energy functionals with a leading length contribution and geometry-dependent corrections \cite{ChuiMoffatt1995,MaggioniRicca2009,Saffman1992}.

This note does not claim that the electron is a vortex.  Its purpose is narrower: to audit a family of conditional normalization formulas, expose their complete dependency structure, and separate exact algebra from physical identification.  In particular, the symbols called ``densities'' below are initially only coefficients in specified energy closures.  They must not be identified with a material substrate density, the canonical effective response density, a local resolved-tube density, or one another without an additional constitutive theorem.

A companion manuscript \cite{Iskandarani2026} studies the numerical relation
\[
\frac{8\pi}{3}\mathcal L_D^{(3_1)}
\simeq \alpha^{-1},
\]
where \(\mathcal L_D^{(3_1)}\) is a diameter-normalized numerical estimate for the ideal trefoil.  That work treats the relation as an independent geometric candidate with a nonzero residual, not as a completed derivation of the fine-structure constant.  The present note asks what follows when this candidate is inserted into several rest-energy bookkeeping closures.

The main results are as follows.

First, the horn-envelope and filament closures imply an exact ratio once their numerical coefficients are frozen.  The ratio becomes \(1/(2\alpha)\) on the calibrated electron chain, but this is a conditional normalization identity, not a universal volume law.

Second, three apparently different numerical comparisons are algebraically identical and therefore constitute only one datum.

Third, the comparison scale \(\GammaC=h/m_e\) eliminates \(\alpha\) from one closure and generates a natural electron-anchored density scale.  This does not prove that the substrate carries that circulation or density.

Fourth, the physically load-bearing questions are not numerical.  They are whether the dynamical energy selects the ideal-rope trefoil, what coefficient replaces the regularized Biot--Savart logarithm, which density enters the line energy, how circulation is selected, and how a defect-sector response matches a bulk Maxwell action.

\section{Conventions, data, and epistemic status}
\label{sec:conventions}

\subsection{CODATA quantities}

We use the CODATA 2022 recommended values \cite{CODATA2022}:
\begin{align}
  \alpha &= 7.297\,352\,5643\times10^{-3},
  &
  \alpha^{-1}&=137.035\,999\,177,
  \notag\\
  \lC &= \frac{\hbar}{m_ec}
  =3.861\,592\,6744\times10^{-13}\ \mathrm{m},
  &
  \re &= \alpha\lC
  =2.817\,940\,3205\times10^{-15}\ \mathrm{m},
  \notag\\
  \frac{h}{m_e}
  &=7.273\,895\,0934\times10^{-4}\ \mathrm{m^2\,s^{-1}}.
\end{align}
CODATA lists both \(\pi\hbar/m_e=h/(2m_e)\) and
\(2\pi\hbar/m_e=h/m_e\) under ``quantum of circulation.''  In this manuscript
\[
\GammaC:=\frac{h}{m_e}
\]
is called the \emph{electron-anchored Compton circulation scale}.  This terminology is deliberately non-dynamical: no microscopic constituent of the proposed medium has been shown to have mass \(m_e\), and no substrate quantization theorem is assumed.

The classical electron radius is used only as the calibrated combination
\[
\re=\frac{e^2}{4\pi\varepsilon_0m_ec^2}
=\alpha\lC.
\]
It is not a measured geometrical size of the electron.

\subsection{Ropelength convention}

For a closed curve with embedded tube radius \(a\) and diameter \(D=2a\), define
\begin{equation}
  \Rop_a:=\frac{\ell}{a},
  \qquad
  \calL_D:=\frac{\ell}{D}
  =\frac12\Rop_a.
  \label{eq:rope-conventions}
\end{equation}
The high-resolution numerical trefoil reported by Przybyl and Pieranski
\cite{PrzybylPieranski2014} has
\begin{equation}
  \Rop_a^{(3_1)}=32.742\,934\,477,
  \qquad
  \calL_D^{(3_1)}=16.371\,467\,2385.
  \label{eq:ropelength}
\end{equation}
These are high-resolution numerical values for a candidate tight configuration, not a rigorous closed-form value of the global minimizer.  Existence and \(C^{1,1}\) regularity of ropelength minimizers are established mathematically \cite{CantarellaKusnerSullivan2002}.

\subsection{Geometric candidate and its single residual}

Define
\begin{equation}
  \alpha_{\rm geom}
  :=\frac{3}{8\pi\calL_D^{(3_1)}}
  =7.291\,112\,3713\times10^{-3},
  \qquad
  \alpha_{\rm geom}^{-1}
  =137.153\,283\,2132.
  \label{eq:alphageom}
\end{equation}
The relative residual is
\begin{equation}
  \delta
  :=\frac{\alpha_{\rm geom}}{\alpha}-1
  =-8.55131\times10^{-4}
  =-855.13\ \mathrm{ppm}.
  \label{eq:residual}
\end{equation}
No \(r_e\), \(r_c\), or calibrated swirl velocity enters the evaluation of
\cref{eq:alphageom}.  The angular coefficient \(8\pi/3\), however, remains a declared normalization whose dynamical origin is an open question \cite{Iskandarani2026}.

For later algebra it is useful to define the \emph{comparison scales}
\begin{align}
  \ell_* &:=\frac{3}{8\pi}\lC,
  \label{eq:ellstar}\\
  D_{\rm geom}
  &:=\frac{\ell_*}{\calL_D^{(3_1)}}
  =\alpha_{\rm geom}\lC.
  \label{eq:Dgeom}
\end{align}
These are reparametrizations of the same candidate relation.  They are not additional predictions.

\subsection{Gate G0: geometric and dynamical minimizers}
\label{subsec:G0}

Ropelength minimization and field-energy minimization are different variational problems.  The use of \(\calL_D^{(3_1)}\) as a physical electron input is justified only if a resolved SST energy functional has a controlled regime in which
\begin{equation}
  E_{\rm SST}[\gamma,a_{\rm core},\ldots]
  =
  T_{\rm eff}\operatorname{Len}(\gamma)
  +E_{\rm exclusion}[\gamma,a_{\rm core}]
  +o(1),
  \label{eq:hard-tube-limit}
\end{equation}
with fixed effective thickness and hard self-avoidance, so that its minimizer converges to a ropelength minimizer.  Soft Hopf-soliton models such as the Faddeev--Skyrme model minimize a full field energy at fixed Hopf charge and need not select the ideal-rope representative \cite{Sutcliffe2007}.  Thus:

\begin{quote}
\textbf{[GATE G0].}
Before \(\calL_D^{(3_1)}\) is used as a physical input, derive or numerically establish a ropelength-dominated limit of the actual SST energy functional.
\end{quote}

\section{Three declared normalization closures}
\label{sec:closures}

Consider a closed filament with centerline length \(\ell\), characteristic radius \(r\), characteristic tangential speed \(v\), and circulation
\begin{equation}
  \Gamma=2\pi rv.
  \label{eq:circulation}
\end{equation}
A resolved vortex energy is generally nonlocal and core dependent.  For a slender filament it contains a logarithmic contribution of the schematic form
\begin{equation}
  E
  \sim
  \frac{\rho\Gamma^2\ell}{4\pi}
  \left[
    \ln\!\left(\frac{R_{\rm out}}{a_{\rm core}}\right)
    +C_{\rm profile}
  \right]
  +E_{\rm nonlocal},
  \label{eq:standard-line-energy}
\end{equation}
where \(C_{\rm profile}\), the outer scale, and nonlocal terms depend on the core model and global geometry \cite{Saffman1992}.  The closures below replace this structure by fixed coefficients.  They are definitions of bookkeeping branches, not derived hydrodynamic theorems.

\begin{definition}[H: horn-envelope normalization]
\label{def:H}
The electron rest energy is represented by
\begin{equation}
  m_ec^2
  =
  C_H\,\rhohorn\,v^2r^3,
  \qquad
  C_H:=2\pi.
  \label{eq:H}
\end{equation}
The coefficient \(\rhohorn\) is a horn-envelope normalization density.  It is not a local resolved-tube density.
\end{definition}

\begin{definition}[F: filament line-energy normalization]
\label{def:F}
The leading line-energy bookkeeping is
\begin{equation}
  m_ec^2
  =
  C_F\,\rhofil\,\Gamma^2\ell,
  \qquad
  C_F:=\frac13.
  \label{eq:F}
\end{equation}
The coefficient \(\rhofil\) is an effective filament energy coefficient until a constitutive theorem identifies it with a physical density.
\end{definition}

\begin{definition}[C: Compton-circulation comparison]
\label{def:C}
As in \cref{def:F}, but with the comparison scale
\begin{equation}
  \Gamma=\GammaC:=\frac{h}{m_e}.
  \label{eq:C}
\end{equation}
The resulting coefficient is denoted \(\rhoCeff\).  This branch is not a quantization theorem for the substrate.
\end{definition}

The dimensions are correct:
\[
[\rho\Gamma^2\ell]
=
\mathrm{kg\,m^{-3}}\,
\mathrm{m^4\,s^{-2}}\,
\mathrm m
=
\mathrm J,
\qquad
[\rho v^2r^3]=\mathrm J.
\]

\section{Conditional ratio of horn and filament normalizations}
\label{sec:ratio}

\begin{proposition}[General coefficient dependence]
\label{prop:general-ratio}
For equal total energy, with \(\Gamma=2\pi rv\), the closures
\[
E=C_H\rho_Hv^2r^3,
\qquad
E=C_F\rho_F\Gamma^2\ell
\]
imply
\begin{equation}
  \boxed{
  \frac{\rho_H}{\rho_F}
  =
  \frac{4\pi^2C_F}{C_H}\frac{\ell}{r}
  }.
  \label{eq:general-ratio}
\end{equation}
The result is independent of \(E\) and \(v\), but it depends explicitly on the closure coefficients \(C_H\) and \(C_F\).
\end{proposition}

\begin{proof}
Solving the two closures gives
\[
\rho_H=\frac{E}{C_Hv^2r^3},
\qquad
\rho_F
=\frac{E}{4\pi^2C_Fv^2r^2\ell}.
\]
Their ratio is \cref{eq:general-ratio}.
\end{proof}

For the declared values \(C_H=2\pi\) and \(C_F=1/3\),
\begin{equation}
  \frac{\rhohorn}{\rhofil}
  =
  \frac{2\pi}{3}\frac{\ell}{r}.
  \label{eq:specific-ratio}
\end{equation}
On the calibrated chain
\begin{equation}
  r=\frac{\re}{2}
  =\frac{\alpha}{2}\lC,
  \qquad
  v=\frac{\alpha c}{2},
  \qquad
  \ell=\ell_*=\frac{3}{8\pi}\lC,
  \label{eq:calibrated-chain}
\end{equation}
this becomes
\begin{equation}
  \boxed{
  \frac{\rhohorn}{\rhofil}
  =
  \frac{1}{2\alpha}
  =68.5180
  }.
  \label{eq:ratio-alpha}
\end{equation}

\begin{remark}[What \cref{eq:ratio-alpha} does and does not mean]
Equation~\eqref{eq:ratio-alpha} is an exact algebraic consequence of the stated coefficients and the calibrated scale chain.  It is not an independent derivation of \(\alpha\), because \(r/\lC=\alpha/2\) is already used.  It is also not a universal geometric volume theorem.

Indeed, the actual Euclidean volumes of a sphere of radius \(r\) and a circular tube of minor radius \(r\) and centerline length \(\ell\) are
\[
V_{\rm sphere}=\frac{4\pi}{3}r^3,
\qquad
V_{\rm tube}^{\rm geom}=\pi r^2\ell,
\]
whose ratio is
\[
\frac{V_{\rm tube}^{\rm geom}}{V_{\rm sphere}}
=\frac34\frac{\ell}{r},
\]
not \((2\pi/3)(\ell/r)\).  The factors \(2\pi r^3\) and
\((4\pi^2/3)r^2\ell\) are effective support factors generated by the two energy normalizations.  They may be useful bookkeeping devices, but they must not be relabelled as literal sphere and tube volumes.
\end{remark}

\begin{remark}[Density non-identification]
The identity compares \(\rhohorn\) and \(\rhofil\) as coefficients in two closures.  It does not establish
\[
\rhofil=\rho_{\!f},
\qquad
\rhofil=\rho_{\rm substrate},
\qquad
\rhofil=\rho_{\rm local\ tube},
\]
or any other density identification.  Such a statement requires a separate constitutive matching theorem.
\end{remark}

Numerically, the three branches give
\begin{equation}
  \rhohorn=3.8934\times10^{18},
  \qquad
  \rhofil=5.6824\times10^{16},
  \qquad
  \rhoCeff=1.0071\times10^{7}
  \quad
  (\mathrm{kg\,m^{-3}}),
  \label{eq:numerical-coefficients}
\end{equation}
but all three numbers are closure outputs built from measured electron quantities and frozen normalization coefficients.

\section{One residual in three equivalent representations}
\label{sec:residual-three}

\begin{proposition}[Residual equivalence]
\label{prop:residual}
With \(\delta\) defined in \cref{eq:residual},
\begin{equation}
  \frac{\alpha_{\rm geom}}{\alpha}-1
  =
  \frac{D_{\rm geom}}{\re}-1
  =
  \frac{3/(4\pi\alpha)}{\Rop_a^{(3_1)}}-1
  =
  \delta.
  \label{eq:three-faces}
\end{equation}
\end{proposition}

\begin{proof}
Because \(D_{\rm geom}=\alpha_{\rm geom}\lC\) and
\(\re=\alpha\lC\),
\[
\frac{D_{\rm geom}}{\re}
=
\frac{\alpha_{\rm geom}}{\alpha}.
\]
Also,
\[
\frac{3/(4\pi\alpha)}{\Rop_a^{(3_1)}}
=
\frac{3}{8\pi\alpha\calL_D^{(3_1)}}
=
\frac{\alpha_{\rm geom}}{\alpha}.
\]
\end{proof}

Thus the agreement of any one representation cannot be used as independent corroboration of the other two.

\subsection{Electromagnetic self-energy matching}

A statement such as \(D=\re\) is not a geometry theorem.  The electromagnetic self-energy of a finite charge distribution has the general form
\begin{equation}
  E_{\rm EM}
  =
  C_{\rm EM}
  \frac{e^2}{4\pi\varepsilon_0D},
  \label{eq:EM-shape}
\end{equation}
where \(C_{\rm EM}\) depends on the charge profile, topology, screening, and boundary conditions.  Requiring \(E_{\rm EM}=m_ec^2\) yields
\begin{equation}
  D=C_{\rm EM}\re.
  \label{eq:shape-match}
\end{equation}
Therefore \(D=\re\) follows only for the special convention \(C_{\rm EM}=1\).  For a knotted or toroidal charge distribution, \(C_{\rm EM}\) must be computed.

It is therefore more accurate to call
\[
E_{\rm hydro}=E_{\rm EM}=m_ec^2
\]
an \emph{equal-self-energy matching ansatz}, not equipartition.

\section{The Compton-circulation comparison branch}
\label{sec:compton-branch}

\subsection{An \texorpdfstring{\(\alpha\)-free}{alpha-free}, electron-anchored density scale}

Using \(\GammaC=h/m_e=2\pi\lC c\) in the filament closure and
\(\ell_*=(3/8\pi)\lC\) gives
\begin{align}
  \rhoCeff
  &=
  \frac{3m_ec^2}{\GammaC^2\ell_*}
  \notag\\
  &=
  \frac{8\pi m_ec^2}{\GammaC^2\lC}
  =
  \boxed{
  \frac{2}{\pi}\frac{m_e}{\lC^3}
  }
  =
  1.0071\times10^7\ \mathrm{kg\,m^{-3}}.
  \label{eq:rhoC}
\end{align}
This expression is independent of \(\alpha\), but it is not parameter free in the predictive sense: it uses the measured \(m_e\), \(\hbar\), \(c\), the geometric coefficient \(3/(8\pi)\), and the line-energy coefficient \(1/3\).  It is an electron-anchored density scale, not a derived substrate density.

Within this frozen closure, \(\alpha\) cancels and the algebra determines a density coefficient rather than the coupling.  This is a structural statement about the closure, not a no-go theorem for every possible SST derivation of \(\alpha\).

\subsection{\texorpdfstring{Formal speed relative to \(c\)}{Formal speed relative to c}}

For a line vortex with \(\GammaC=2\pi\lC c\),
\begin{equation}
  v_\theta(r)
  =
  \frac{\GammaC}{2\pi r}
  =
  c\frac{\lC}{r}.
  \label{eq:vGammaC}
\end{equation}
At \(a=D_{\rm geom}/2\), the formal ratio is
\[
\frac{v_\theta(a)}{c}
=
\frac{2}{\alpha_{\rm geom}}
\simeq274.3,
\]
and at \(r=2\ell_*\),
\[
\frac{v_\theta(2\ell_*)}{c}
=
\frac{4\pi}{3}
\simeq4.19.
\]

These numbers do not contradict Galilean incompressible Euler theory, in which \(c\) is not a fundamental speed limit and the pressure constraint is instantaneous.  They do obstruct a literal identification of this azimuthal transport speed with a physical matter velocity in a sector obeying universal operational Lorentz invariance with limiting speed \(c\).  A viable model must therefore supply a separate emergent-causality map, a finite-core modification of \cref{eq:vGammaC}, or a different circulation sector \cite{BarceloLiberatiVisser2026}.

\subsection{Calibrated circulation ratio}

On the calibrated horn chain,
\begin{equation}
  r=\frac{\alpha}{2}\lC,
  \qquad
  v=\frac{\alpha}{2}c,
\end{equation}
so
\begin{equation}
  \Gammaeff
  =
  2\pi rv
  =
  \frac{\alpha^2}{4}\frac{h}{m_e},
  \qquad
  \boxed{
  \frac{\Gammaeff}{\GammaC}
  =
  \frac{\alpha^2}{4}
  }.
  \label{eq:gamma-ratio}
\end{equation}
This ratio is exact, but it is not an independent prediction: both \(r\) and \(v\) are \(\alpha\)-calibrated.  It becomes a physical ``circulation deficit'' only if a separate theorem establishes that \(\GammaC=h/m_e\) is the circulation quantum that the same defect ought to carry.  Until then, \cref{eq:gamma-ratio} is a transparent comparison of two normalization scales.

The corresponding coefficient ratio is
\begin{equation}
  \frac{\rhofil}{\rhoCeff}
  =
  \left(\frac{\GammaC}{\Gammaeff}\right)^2
  =
  \frac{16}{\alpha^4}.
  \label{eq:rho-ratio-gamma}
\end{equation}
Again, this is power counting in a calibrated chain, not evidence for a microscopic suppression mechanism.

\begin{table}[t]
\caption{Conditional closure outputs on the chain \cref{eq:calibrated-chain}.  The entries are normalization coefficients until constitutive identifications are derived.}
\label{tab:closures}
\begin{ruledtabular}
\begin{tabular}{llll}
Branch & defining circulation & coefficient (kg\,m\(^{-3}\)) & status\\
\hline
H & \(\Gammaeff\) implicit & \(3.8934\times10^{18}\) & horn-envelope calibration\\
F & \(\Gammaeff=(\alpha^2/4)h/m_e\) & \(5.6824\times10^{16}\) & filament closure, \(\alpha\)-contaminated\\
C & \(\GammaC=h/m_e\) & \(1.0071\times10^{7}\) & \(\alpha\)-free comparison scale\\
\end{tabular}
\end{ruledtabular}
\end{table}

\section{Meaning of the filament coefficient}
\label{sec:line-coefficient}

Comparing the declared filament closure
\[
E=\frac13\rho\Gamma^2\ell
\]
with the leading hollow-core expression
\[
E
\simeq
\frac{\rho\Gamma^2\ell}{4\pi}
\left[
\ln\!\left(\frac{R_{\rm out}}{a_{\rm core}}\right)
+C_{\rm profile}
\right]
\]
shows that the effective condition is
\begin{equation}
  \ln\!\left(\frac{R_{\rm out}}{a_{\rm core}}\right)
  +C_{\rm profile}
  +C_{\rm nonlocal}
  =
  \frac{4\pi}{3}.
  \label{eq:effective-log}
\end{equation}
If all additive profile and nonlocal constants were set to zero, one would obtain \(R_{\rm out}/a_{\rm core}=e^{4\pi/3}=65.94\).  That number is not a computed physical cutoff: an \(O(1)\) change in the additive constants changes the inferred ratio multiplicatively and is vastly larger than the \(855\)-ppm residual.

Hence:

\begin{quote}
\textbf{[GATE G1].}
Compute the finite-core and nonlocal coefficient of the actual SST trefoil from a declared velocity profile and regulator.  Do not infer sub-per-mille physics from the provisional value \(C_F=1/3\).
\end{quote}

A second theorem is required:

\begin{quote}
\textbf{[GATE G2].}
Determine whether the coefficient multiplying \(\Gamma^2\ell\) contains the material substrate density, the quasi-static effective response density, a local tube density, a line-inertia coefficient, or another constitutive quantity.
\end{quote}

\section{Maxwell matching and electron constraints}
\label{sec:constraints}

\subsection{Maxwell/source bridge}

The scale relation \(D/\lC\) is not by itself an electromagnetic derivation.  A genuine bridge requires three logically separate results.  In the action convention
\begin{equation}
  \frac{S_{\rm IR}^{(2)}}{\hbar}
  =
  -\frac{\mathscr R_{\rm EM}}{16\pi}
  \int f_{\mu\nu}f^{\mu\nu}\,\dd^4x
  +
  \int j^\mu a_\mu\,\dd^4x,
  \label{eq:IR-action}
\end{equation}
with \(g^2=4\pi\alpha\), the target coefficient is
\begin{equation}
  \mathscr R_{\rm EM}
  =
  \frac{4\pi}{g^2}
  =
  \alpha^{-1}.
  \label{eq:EM-stiffness}
\end{equation}
The three gates are:

\begin{enumerate}[label=\textbf{B2\alph*:},leftmargin=*]
  \item derive a local, gauge-invariant, monometric Maxwell quadratic action in the observable infrared sector;
  \item derive the source normalization that assigns the elementary defect the observed unit charge;
  \item prove that the hydrodynamic/defect response coefficient computed from the vortex sector equals the bulk electromagnetic stiffness \(\mathscr R_{\rm EM}\).
\end{enumerate}

The third step is the load-bearing one.  A photon kinetic coefficient is a property of the bulk response; its equality to a single-defect geometric functional cannot be imposed by notation.

\subsection{Magnetic-moment constraint}

The electron magnetic moment is measured with \(0.13\) parts-per-trillion precision \cite{Fan2023}.  Translating that precision into a compositeness length is model dependent.  For a chirality-protected dipole correction,
\begin{equation}
  \Delta a_e
  \sim
  \frac{m_e^2}{\Lambda^2},
\end{equation}
an illustrative allowance \(|\Delta a_e|\lesssim10^{-12}\) gives
\begin{equation}
  \Lambda\gtrsim511\ \mathrm{GeV},
  \qquad
  \frac{\hbar c}{\Lambda}
  \lesssim3.9\times10^{-19}\ \mathrm m.
  \label{eq:gminus2-bound}
\end{equation}
Without chirality protection the scaling may be linear and the bound much stronger \cite{BrodskyDrell1980,GiudiceParadisiPassera2012}.

Thus a femtometre-scale structure that carries charge, magnetic moment, or an unsuppressed electromagnetic form factor is excluded.  This does not directly exclude an electromagnetically inert mass-normalization structure.  Any such separation between a mass-bearing filament and a pointlike electromagnetic sector must itself be derived and tested.

\section{Preregistered falsification gates}
\label{sec:programme}

The present article supports the following gate structure.

\begin{description}[leftmargin=2.3cm,style=nextline]
  \item[\textbf{G0: dynamics}]
  Show that the resolved SST energy selects the ideal-rope trefoil, or replace \(\calL_D^{(3_1)}\) by the dynamically selected value.

  \item[\textbf{G1: line energy}]
  Compute \(C_F\), the core-profile contribution, and the nonlocal Biot--Savart contribution without fitting to \(\alpha\).

  \item[\textbf{G2: density sector}]
  Derive which material, response, local-tube, or line-inertia coefficient enters the filament energy.

  \item[\textbf{G3: circulation}]
  Derive the physical circulation of the defect.  In particular, establish or reject \(\GammaC=h/m_e\) without using \(r_e\), \(r_c\), or \(\alpha c\) as hidden calibration inputs.

  \item[\textbf{G4: electromagnetism}]
  Derive B2a--B2c and compute the self-energy shape factor \(C_{\rm EM}\).

  \item[\textbf{G5: out-of-sample test}]
  After all coefficients and profiles are frozen, compare the resulting dimensionless coupling with \(\alpha\).  The measured \(\alpha\), \(r_e\), and any transitive proxies are forbidden during the fit or selector stage.
\end{description}

\begin{table}[t]
\caption{Epistemic status of the central statements.}
\label{tab:status}
\begin{ruledtabular}
\begin{tabular}{ll}
Statement & status\\
\hline
\(\alpha_{\rm geom}^{-1}=(8\pi/3)\calL_D^{(3_1)}\) & independent geometric candidate\\
Residual equivalence, \cref{eq:three-faces} & algebraic identity\\
\(\rhohorn/\rhofil=1/(2\alpha)\) & calibrated conditional identity\\
Literal volume interpretation of that ratio & not derived\\
\(\rhoCeff=(2/\pi)m_e/\lC^3\) & \(\alpha\)-free electron-anchored scale\\
\(\rhoCeff=\rho_{\!f}\) or substrate density & open / not identified\\
\(\Gammaeff/\GammaC=\alpha^2/4\) & calibrated comparison\\
\(\GammaC\) as substrate circulation quantum & open\\
\(D=\re\) & special case \(C_{\rm EM}=1\), not a theorem\\
Ideal trefoil as physical minimizer & open Gate G0\\
\end{tabular}
\end{ruledtabular}
\end{table}

\section{Conclusion}
\label{sec:conclusion}

The displayed normalization formulas are algebraically coherent once their coefficients and calibration ancestry are stated explicitly.  Their interpretation is substantially more limited than the raw numerical coincidences suggest.

The factor \(1/(2\alpha)\) is exact under the chosen horn and filament normalizations, but it is not a universal ratio of geometric volumes and is not independent of the calibrated relation \(r/\lC=\alpha/2\).  The three \(855\)-ppm comparisons are one equation in three forms.  The Compton-circulation branch produces an \(\alpha\)-free density scale, but neither its circulation nor its density has been derived as a substrate property.  The factor \(\alpha^2/4\) is likewise a transparent calibrated ratio rather than a demonstrated dynamical suppression.

Consequently, the model cannot presently be described as containing exactly one physical free parameter.  Once the provisional conventions are unfrozen, at least the dynamical knot representative, the line-energy coefficient, the relevant density sector, the circulation map, the electromagnetic stiffness match, and the electrostatic shape factor remain open.

This is nevertheless a useful result.  The problem is now decomposed into independent theorem and falsification gates, and hidden reuse of \(\alpha\) is visible by inspection.  Progress requires a resolved core calculation and an action-level electromagnetic match, not further manipulation of the same calibrated identities.

\begin{acknowledgments}
The author thanks the maintainers of the high-resolution ideal-knot data and the CODATA adjustment for making precise reference values available.
\end{acknowledgments}

\appendix

\section{Power counting in \texorpdfstring{\(\varepsilon=\alpha/2\)}{epsilon}}
\label{app:powers}

Let
\[
\varepsilon:=\frac{\alpha}{2},
\qquad
r=\varepsilon\lC,
\qquad
v=\varepsilon c.
\]
Then
\begin{equation}
  \Gammaeff
  =
  2\pi\varepsilon^2\lC c,
  \qquad
  \GammaC
  =
  2\pi\lC c,
  \qquad
  \ell_*=\frac{3}{8\pi}\lC.
\end{equation}
The three closure coefficients are
\begin{equation}
  \rhohorn
  =
  \frac{m_e}{2\pi\varepsilon^5\lC^3},
  \qquad
  \rhofil
  =
  \frac{2m_e}{\pi\varepsilon^4\lC^3},
  \qquad
  \rhoCeff
  =
  \frac{2m_e}{\pi\lC^3}.
\end{equation}
Therefore
\begin{equation}
  \frac{\rhohorn}{\rhofil}
  =
  \frac{1}{4\varepsilon}
  =
  \frac{1}{2\alpha},
  \qquad
  \frac{\rhofil}{\rhoCeff}
  =
  \varepsilon^{-4}
  =
  \frac{16}{\alpha^4}.
\end{equation}
The first ratio contains one inverse power of \(\alpha\) because the two closures scale as \(\varepsilon^{-5}\) and \(\varepsilon^{-4}\).  The second is the square of the calibrated circulation ratio.  These equations are provenance diagnostics, not independent predictions of \(\alpha\).


\begin{thebibliography}{99}

\bibitem{Thomson1867}
W.~Thomson (Lord Kelvin),
``On vortex atoms,''
\emph{Phil. Mag.} \textbf{34}, 15--24 (1867).

\bibitem{Onsager1949}
L.~Onsager,
``Statistical hydrodynamics,''
\emph{Nuovo Cimento Suppl.} \textbf{6}, 279--287 (1949),
\href{https://doi.org/10.1007/BF02780991}{doi:10.1007/BF02780991}.

\bibitem{Feynman1955}
R.~P.~Feynman,
``Application of quantum mechanics to liquid helium,''
in \emph{Progress in Low Temperature Physics}, Vol.~1, edited by C.~J.~Gorter
(North-Holland, Amsterdam, 1955), pp.~17--53,
\href{https://doi.org/10.1016/S0079-6417(08)60077-3}{doi:10.1016/S0079-6417(08)60077-3}.

\bibitem{Vinen1961}
W.~F.~Vinen,
``The detection of single quanta of circulation in liquid helium II,''
\emph{Proc. R. Soc. Lond. A} \textbf{260}, 218--236 (1961),
\href{https://doi.org/10.1098/rspa.1961.0029}{doi:10.1098/rspa.1961.0029}.

\bibitem{Donnelly1991}
R.~J.~Donnelly,
\emph{Quantized Vortices in Helium II}
(Cambridge University Press, Cambridge, 1991).

\bibitem{KlecknerIrvine2013}
D.~Kleckner and W.~T.~M.~Irvine,
``Creation and dynamics of knotted vortices,''
\emph{Nat. Phys.} \textbf{9}, 253--258 (2013),
\href{https://doi.org/10.1038/nphys2560}{doi:10.1038/nphys2560}.

\bibitem{BuniyKephart2005}
R.~V.~Buniy and T.~W.~Kephart,
``Glueballs and the universal energy spectrum of tight knots and links,''
\emph{Int. J. Mod. Phys. A} \textbf{20}, 1252--1259 (2005),
\href{https://doi.org/10.1142/S0217751X05024146}{doi:10.1142/S0217751X05024146}.

\bibitem{BuniyEtAl2014}
R.~V.~Buniy, J.~Cantarella, T.~W.~Kephart, and E.~J.~Rawdon,
``Tight knot spectrum in QCD,''
\emph{Phys. Rev. D} \textbf{89}, 054513 (2014),
\href{https://doi.org/10.1103/PhysRevD.89.054513}{doi:10.1103/PhysRevD.89.054513}.

\bibitem{ChuiMoffatt1995}
A.~Y.~K.~Chui and H.~K.~Moffatt,
``The energy and helicity of knotted magnetic flux tubes,''
\emph{Proc. R. Soc. Lond. A} \textbf{451}, 609--629 (1995),
\href{https://doi.org/10.1098/rspa.1995.0146}{doi:10.1098/rspa.1995.0146}.

\bibitem{MaggioniRicca2009}
F.~Maggioni and R.~L.~Ricca,
``On the groundstate energy of tight knots,''
\emph{Proc. R. Soc. A} \textbf{465}, 2761--2783 (2009),
\href{https://doi.org/10.1098/rspa.2008.0536}{doi:10.1098/rspa.2008.0536}.

\bibitem{Iskandarani2026}
O.~Iskandarani,
``Angular-projector normalization and finite-core energetics of tight knotted vortex tubes,''
preprint, manuscript version 0.2.0 (2026).

\bibitem{CODATA2022}
P.~J.~Mohr, D.~B.~Newell, B.~N.~Taylor, and E.~Tiesinga,
``CODATA recommended values of the fundamental physical constants: 2022,''
\emph{Rev. Mod. Phys.} \textbf{97}, 025002 (2025),
\href{https://doi.org/10.1103/RevModPhys.97.025002}{doi:10.1103/RevModPhys.97.025002}.

\bibitem{PrzybylPieranski2014}
S.~Przybyl and P.~Pieranski,
``High resolution portrait of the ideal trefoil knot,''
\emph{J. Phys. A: Math. Theor.} \textbf{47}, 285201 (2014),
\href{https://doi.org/10.1088/1751-8113/47/28/285201}{doi:10.1088/1751-8113/47/28/285201}.

\bibitem{CantarellaKusnerSullivan2002}
J.~Cantarella, R.~B.~Kusner, and J.~M.~Sullivan,
``On the minimum ropelength of knots and links,''
\emph{Invent. Math.} \textbf{150}, 257--286 (2002),
\href{https://doi.org/10.1007/s00222-002-0234-y}{doi:10.1007/s00222-002-0234-y}.

\bibitem{Saffman1992}
P.~G.~Saffman,
\emph{Vortex Dynamics}
(Cambridge University Press, Cambridge, 1992),
\href{https://doi.org/10.1017/CBO9780511624063}{doi:10.1017/CBO9780511624063}.

\bibitem{Ohanian1986}
H.~C.~Ohanian,
``What is spin?,''
\emph{Am. J. Phys.} \textbf{54}, 500--505 (1986),
\href{https://doi.org/10.1119/1.14580}{doi:10.1119/1.14580}.

\bibitem{BarceloLiberatiVisser2026}
C.~Barcel\'o, S.~Liberati, and M.~Visser,
``Analogue gravity,''
\emph{Living Rev. Relativ.} \textbf{29}, 2 (2026),
\href{https://doi.org/10.1007/s41114-026-00064-9}{doi:10.1007/s41114-026-00064-9}.

\bibitem{Fan2023}
X.~Fan, T.~G.~Myers, B.~A.~D.~Sukra, and G.~Gabrielse,
``Measurement of the electron magnetic moment,''
\emph{Phys. Rev. Lett.} \textbf{130}, 071801 (2023),
\href{https://doi.org/10.1103/PhysRevLett.130.071801}{doi:10.1103/PhysRevLett.130.071801}.

\bibitem{BrodskyDrell1980}
S.~J.~Brodsky and S.~D.~Drell,
``Anomalous magnetic moment and limits on fermion substructure,''
\emph{Phys. Rev. D} \textbf{22}, 2236--2243 (1980),
\href{https://doi.org/10.1103/PhysRevD.22.2236}{doi:10.1103/PhysRevD.22.2236}.

\bibitem{Sutcliffe2007}
P.~Sutcliffe,
``Knots in the Skyrme--Faddeev model,''
\emph{Proc. R. Soc. A} \textbf{463}, 3001--3020 (2007),
\href{https://doi.org/10.1098/rspa.2007.0038}{doi:10.1098/rspa.2007.0038}.


\bibitem{GiudiceParadisiPassera2012}
G.~F.~Giudice, P.~Paradisi, and M.~Passera,
``Testing new physics with the electron \(g-2\),''
\emph{J. High Energy Phys.} \textbf{11}, 113 (2012),
\href{https://doi.org/10.1007/JHEP11(2012)113}{doi:10.1007/JHEP11(2012)113}.

\end{thebibliography}

\end{document}
